% !TeX root = ../papers/nc-p3.tex
\section{階数15の通常倍率の像を補う四つの正則窓}
\label{sec:rank-fifteen-formal-boundary}
\begin{theorem}\label{thm:rank-fifteen-formal-boundary}
$y\ne0$, $p^2=x^3-27y$, $q^2=(1-x)^2-4y$ の規格化テンソルの商曲面を考える。
通常倍率の階数15窓の像は、この曲面から次の四点を除いたものにちょうど等しい。
全て $q=0$ 上にあり、$y=(1-x)^2/4$ である。
\begin{enumerate}
\item $x=3,3/4$ の各点（$p=0$）。
\item $U=3\pm2\sqrt2$, $x=(U+3)/4$,
$p=-U^2(U-9)/(U-1)^3$ の各点。
\end{enumerate}
四点の規格化テンソルも標準14641型のKoszul左右AS正則窓を与え、
各始点で四中心四次式の多項式環上自由階数256となる。
最初の二点の最小周期は1、残り二点の最小周期は2である。
\end{theorem}
\begin{proof}
$K=rk$, $H=rm$, $r=s^2$, $w=t^2$ とする。
形式的表示では $U=K^2\ne1$, $V=H^2\ne0$ を常に満たす。
通常条件はこの上で
\[
 K\ne0,\qquad K+H\notin\{0,1,-1\},\qquad K-H\notin\{1,-1\}
\]
に等しい。実際 $2(K+H)=r+1/r$, $2(K-H)=w+1/w$ である。
$w=-r,-1/r$ は $K=0$ に一致し、$w=r,1/r$ は $H=0$ に一致する。
また $z=-1$ は $K+H=0$ を与え、$z=1$ と $z\lambda^2=1$ は各和差が $\pm1$ となる条件である。

$q\ne0$ なら $K\pm H=\pm1$ は $q$ の積表示から起こらない。
$p\ne0$ なら三次式 $v(x-4v)^2-y$ の判別式は $256yp^2\ne0$ である。
その三根のうち $3v^2-xv+y\ne0$ となるものを取ると $U=x/v-3$ と $V=y/v^2$ は異なる。
$p\ne0$ により $K\ne0$ で、符号のどちらでも $K+H\ne0$ となる。
$KH$ の符号で所定の $p$ を、$w$ の逆転で所定の $q$ を選べる。
$p=0$, $q\ne0$ なら $x\ne0$ で $U=9$, $V=16x/3$ を選べばよい。
$K=3$ として $H$ の符号で $K+H\ne0$ にでき、同様に所定の $q$ になる。

$q=0$ で通常なら $2n^2+1-U=0$ が必要である。よって
$V=(U-1)^2/4$, $U=4x-3$, $H=\pm(U-1)/2$ である。
このとき $((K+H)^2-1)((K-H)^2-1)=(U-9)(U-1)^3/16$ なので、
$U=9$ と $K=0$ は除かれる。それ以外では $p$ が $KH$ を一意に指定する。
最後の例外 $K+H=0$ は $U=V$ と $KH=-U$ に相当し、
$(U-1)^2=4U$ を解いて記載した二点を得る。全ての除外条件を避けたときは
二次式から $r,w,n$ を復元できるので十分性も従う。

四点では形式的表示 $K=0,3,1+\sqrt2,1-\sqrt2$ を順に用いる。
初めの二つで $H=(U-1)/2$、後の二つで $H=-K$ とし、$n^2=(U-1)/2$ を取る。
\path{computations/verify_rank_fifteen_formal_boundary.py} は17の上で二平方根の符号と
$r,w,s,t,n$ の根を選び、四テンソルを有限体で直接作る。選んだ還元は逐次モニックな
方程式の根として標数零へ持ち上がり、上記の四点を尽くす。
各巡回隣接像の階数は6で、その直交補空間を二次関係とする双対の経路次元は
全始点で長さ0から5まで $(1,4,6,4,1,0)$ となる。
元の経路次元と四次中心の文字次元を確認し、四中心商も全始点で長さ13まで所定の次元列を得る。

この場合の持上げの下界は形式的パラメータ空間 $X$ の既約性から得る。
その稠密な通常窓で隣接階数6なので、剰余で階数6なら標数零の点とその近傍で関係は定階数である。
経路と双対経路の次元は正則一般点からの下界に等しく、四次中心も経路基底を取って同様に持ち上がる。
反対階数15は十五係数の非消滅から従い、各一文字切断の階数も4である
（一文字の台が三次元なら反対階数は高々12）。四次双対の双対は元のテンソル直線なので、
一文字と隣接二文字の切断階数から双対の全次数の完全対合を得る。
四中心商の上界と元の経路の全次数の下界を合わせた自由性・Koszul完全性の議論
（定理\ref{thm:klein-rational-period-four}）が適用できる。
また四次テンソルを関係から回収できるので、三つの縮約不変量はこの四窓でも添字保持不変であり、
通常倍率の窓とも互いにも混同されない。

形式的な $t$ の逆転は語反転である。実際Pauliテンソルの二因子を転置して交換すると、
ラベル $(a,b)$ は $(a,\beta b)$ に移り、符号は $(-1)^{e((a,b),(a,b))}$ になる。
コードはその係数配列が $C,D$ を $-D,-C$ に替えた配列と一致することを確認する。
$s,t$ の交換は既に第一因子の転置と共役である。従って一位置移動は
$(p,q)\mapsto(-p,-q)$ に対応し、二位置移動は元の192軌道に戻る。
軌道の分離と上の縮約不変性から記載した最小周期が従う。
\end{proof}

\begin{corollary}
上の形式的な階数15の族では $q=0$ の曲線の全点が正則である。
その正規化は $j\ne\pm1$ の直線で、
$x=(j^2+3)/4$, $p=j(j^2-9)/8$、$j=\pm3$ が同じ節点に移る。
\end{corollary}
\begin{proof}
通常点では $q=0$ は正則な分岐であり、残る四点は定理で認証した。
$p^2=(x-3)^2(4x-3)/4$ に $j=2p/(x-3)$ を代入すれば表示を得る。
\end{proof}
